\documentclass[11pt]{article}
\usepackage[margin=1in]{geometry}
\usepackage{graphicx}

\title{CSE 351: Artificial Intelligence -- Assignment 02}
\author{Student ID 1, Student ID 2}
\date{Fall 2026}

\begin{document}
\maketitle

\section*{Question 1: Neural Networks as Universal Approximators}

\subsection*{Baseline Model Results}
A Multi-Layer Feedforward Neural Network (MFNN) was trained using backpropagation for 5000 epochs with 1 hidden layer (40 neurons), \texttt{tansig} activation, learning rate of 0.05, and initial bias of 0.5 to approximate $f(x, y) = e^{-\frac{x^2+y^2}{8}}\cos(\sqrt{x^2+y^2}+c)$ with $c = 1 \times 10^{-3}$.

\begin{table}[h!]
\centering
\begin{tabular}{|l|c|}
\hline
\textbf{Metric} & \textbf{Value} \tabularnewline \hline
Initial Cost $J$ (Epoch 0) & 0.976229 \tabularnewline \hline
Final Cost $J$ (Epoch 5000) & 0.000185 \tabularnewline \hline
Training MSE & 0.000370 \tabularnewline \hline
Testing MSE & 0.000394 \tabularnewline \hline
\end{tabular}
\caption{Baseline MFNN Training and Testing Results}
\end{table}

% Uncomment the lines below to include your saved plots:
% \begin{center}
% \includegraphics[width=0.45\textwidth]{cost_vs_epoch.png}
% \includegraphics[width=0.45\textwidth]{pred_vs_true.png}
% \end{center}

\subsection*{Fine-Tuning and Hyperparameter Findings}

\begin{table}[h!]
\centering
\begin{tabular}{|l|l|c|c|}
\hline
\textbf{Hyperparameter} & \textbf{Configuration} & \textbf{Train MSE} & \textbf{Test MSE} \tabularnewline \hline
Learning Rate & lr = 0.01 & 0.003373 & 0.003497 \tabularnewline
              & lr = 0.05 (Baseline) & 0.000370 & 0.000394 \tabularnewline
              & lr = 0.1 & 0.000116 & 0.000124 \tabularnewline
              & lr = 0.3 & 0.000127 & 0.000150 \tabularnewline \hline
Number of Layers & 1 Layer: [40] (Baseline) & 0.000370 & 0.000394 \tabularnewline
                 & 2 Layers: [40, 20] & 0.000051 & 0.000060 \tabularnewline
                 & 3 Layers: [40, 20, 10] & 0.000038 & 0.000049 \tabularnewline \hline
Activation Function & tansig (Baseline) & 0.000370 & 0.000394 \tabularnewline
                    & logsig & 0.004096 & 0.004191 \tabularnewline
                    & relu & 0.000284 & 0.000462 \tabularnewline \hline
Initial Bias & bias = 0.0 & 0.001614 & 0.001686 \tabularnewline
             & bias = 0.5 (Baseline) & 0.000370 & 0.000394 \tabularnewline
             & bias = 1.0 & 0.000063 & 0.000068 \tabularnewline
             & bias = 2.0 & 0.000073 & 0.000090 \tabularnewline \hline
\end{tabular}
\caption{Effect of Hyperparameters on Train and Test MSE}
\end{table}

\noindent \textbf{Findings:}
\begin{itemize}
    \item \textbf{Number of Layers:} A single hidden layer of 40 neurons accurately approximates the function, while adding 2 and 3 hidden layers further reduces Test MSE from 0.000394 down to 0.000049 due to higher model capacity.
    \item \textbf{Activation Function:} \texttt{tansig} achieves the best generalization (Test MSE: 0.000394) because it is smooth and zero-centered. \texttt{logsig} converges much slower (Test MSE: 0.004191) due to smaller gradient derivatives, while \texttt{relu} produces piecewise-linear segments that slightly overfit the training points (lower Train MSE of 0.000284, but worse Test MSE of 0.000462).
    \item \textbf{Biases:} Setting initial bias to 0.0 gives higher error (Test MSE: 0.001686) because all hidden neurons output zero at the origin $(0,0)$ where the target function peaks at 1.0. Non-zero biases (0.5 and 1.0) break symmetry at the origin and lower Test MSE to 0.000068, though pushing bias too high (2.0) begins to saturate the neurons.
    \item \textbf{Learning Rates:} A small learning rate (0.01) converges too slowly within 5000 epochs (Test MSE: 0.003497). Increasing the rate to 0.1 yields the lowest error (Test MSE: 0.000124), whereas 0.3 introduces gradient oscillations that slightly degrade final accuracy (Test MSE: 0.000150).
\end{itemize}

\section*{Question 2: Decision Tree Induction}

\begin{table}[h!]
\centering
\begin{tabular}{|l|c|c|c|c|c|c|}
\hline
\textbf{Dataset} & $a_1$ & $a_2$ & $a_3$ & $a_4$ & $a_5$ & $a_6$ \tabularnewline \hline
Data1 & TBD & TBD & TBD & TBD & TBD & TBD \tabularnewline \hline
Data2 & TBD & TBD & TBD & TBD & TBD & TBD \tabularnewline \hline
Data3 & TBD & TBD & TBD & TBD & TBD & TBD \tabularnewline \hline
\end{tabular}
\caption{GINI Index of Attributes Across Datasets}
\end{table}

\section*{Question 3: Image Classification for Object Recognition}

\begin{table}[h!]
\centering
\begin{tabular}{|l|c|}
\hline
\textbf{Classifier} & \textbf{Test Accuracy (Percent)} \tabularnewline \hline
(a) Decision Tree & TBD \tabularnewline \hline
(b) Artificial Neural Network & TBD \tabularnewline \hline
(c) Support Vector Machine (SVM) & TBD \tabularnewline \hline
(d) K-Nearest Neighbour (kNN) & TBD \tabularnewline \hline
\end{tabular}
\caption{CIFAR-10 Classification Accuracy}
\end{table}

\end{document}
